\section{Алгоритм построения реферата}
\label{sec:algorithms}

По варианту~22 обрабатываются русский и немецкий тексты двух областей: научные материалы по медицине и критика предметов изобразительного искусства. Методика \mdash sentence extraction и семантическая сеть OSTIS.

Сначала из файла извлекается текст. Первая непустая строка считается заголовком. Язык определяется по доле кириллицы. Дальше документ режется на абзацы по пустой строке и на предложения по точке, вопросу и восклицанию.

В вес не входят числа, стоп-слова и слова чужого алфавита: латиница в русском тексте, кириллица в немецком. Русские словоформы приводятся к нормальной форме словарём pymorphy3, немецкие \mdash стеммером Snowball. Так выполняется требование методички не учитывать латинские вставки, цифры и служебную лексику.

Базовый вес слова считается по модифицированной TF-IDF:

\begin{equation}
w(t, D) = 0{,}5 \left(1 + \frac{\mathrm{tf}(t, D)}{\mathrm{tf}_{\max}(D)}\right) \log\frac{|DB|}{\mathrm{df}(t)}.
\end{equation}

Вес предложения \mdash произведение трёх функций:

\begin{equation}
\mathrm{Posd}(S_i) = 1 - \frac{BD(S_i)}{|D|}, \qquad
\mathrm{Posp}(S_i) = 1 - \frac{BP(S_i)}{|P|},
\end{equation}

\begin{equation}
\mathrm{Score}(S_i) = \sum_{t \in S_i} \mathrm{tf}(t, S_i)\, w(t, D),
\end{equation}

\begin{equation}
W(S_i) = \mathrm{Posd}(S_i) \cdot \mathrm{Posp}(S_i) \cdot \mathrm{Score}(S_i).
\end{equation}

Здесь $|D|$ \mdash число символов документа, $BD(S_i)$ \mdash символы до предложения в документе, $|P|$ \mdash длина абзаца, $BP(S_i)$ \mdash символы до предложения в абзаце. В реферат попадают десять предложений с наибольшим $W$, но в том порядке, в каком они стоят в тексте.

Позиционные множители поднимают начало документа и начало абзаца. Поэтому метод не сводится к «взять первые десять»: позднее предложение с редкими терминами может обойти раннее бытовое. Для сравнения рядом считается базовый реферат \mdash ровно первые десять предложений без весов.

Реферат ключевых слов строится как дерево. Опорные вершины \mdash восемь лексем с наибольшим $\mathrm{tf}\cdot\mathrm{IDF}$. Если пара соседних лексем встретилась не меньше двух раз и одно из слов уже в дереве, словосочетание вешается ребёнком. Это повторяет пример методички: «лазер / синий лазер».

Часть OSTIS \mdash семантическая сеть: узлы-понятия и дуги-отношения. Документ относится к предметной области и связан с ключевыми понятиями. Понятие включает словосочетание, принадлежит области и совстречается с другим понятием, если оба оказались в одном предложении.

Точность и полнота считаются по множеству номеров предложений. Эталон размечен вручную: в него входят тезисы о тяжести, схеме лечения, деэскалации либо о композиции, смещении и отказе сводить работу к лозунгу. Бытовые описания кашля или «просто чёрного пятна» в эталон не входят.
